\documentclass[12pt,a4paper,oneside]{book}

\usepackage{fontspec}
\usepackage{polyglossia}
\setmainlanguage{russian}
\setotherlanguage{english}
\setmainfont{Times New Roman}

\usepackage{amsmath,amssymb,amsthm,mathtools}
\usepackage{geometry}
\usepackage{enumitem}
\usepackage{array}
\usepackage{booktabs}
\usepackage{hyperref}
\usepackage{microtype}

\geometry{margin=2.5cm}
\setlength{\parindent}{1.25cm}
\setlength{\parskip}{0.35em}
\sloppy

\hypersetup{
    colorlinks=true,
    linkcolor=black,
    citecolor=black,
    urlcolor=blue,
    pdftitle={Теория Единого Целого. Полная согласованная версия},
    pdfauthor={Andrey Tikhonov (Тихонов Андрей)}
}

\theoremstyle{definition}
\newtheorem{axiom}{Аксиома}[chapter]
\newtheorem{definition}{Определение}[chapter]
\newtheorem{postulate}{Постулат}[chapter]

\theoremstyle{plain}
\newtheorem{theorem}{Теорема}[chapter]
\newtheorem{lemma}{Лемма}[chapter]
\newtheorem{corollary}{Следствие}[chapter]
\newtheorem{proposition}{Предложение}[chapter]

\theoremstyle{remark}
\newtheorem{remark}{Замечание}[chapter]
\newtheorem{example}{Пример}[chapter]

% --- Первичные символы и структуры ---
\newcommand{\Whole}{U}
\newcommand{\Parts}{\mathcal A}
\newcommand{\SubParts}{\mathcal B}
\newcommand{\Rel}{R}
\newcommand{\Values}{V}
\newcommand{\Vstruct}{\mathbf V}
\newcommand{\State}{S}
\newcommand{\States}{\mathcal S}
\newcommand{\PartOf}{\sqsubseteq}
\newcommand{\PropPart}{\sqsubset}
\newcommand{\comp}{\circ}
\newcommand{\Vord}{\preceq}
\newcommand{\norm}{\rho}
\newcommand{\ide}{e}
\newcommand{\Disc}{\mathsf{Disc}}
\newcommand{\Space}{\mathsf{Space}}
\newcommand{\Time}{\mathsf{Time}}
\newcommand{\Physics}{\mathsf{Physics}}
\newcommand{\Motion}{\mathsf{Motion}}
\newcommand{\Observer}{\mathsf{Obs}}
\newcommand{\Struct}{\operatorname{Struct}}
\newcommand{\Inv}{\operatorname{Inv}}
\newcommand{\Center}{\operatorname{Center}}
\newcommand{\Varr}{\operatorname{Var}}
\newcommand{\Stab}{\sigma}
\newcommand{\Cost}{\mathcal C}

\title{\Huge\textbf{Теория Единого Целого}\\[0.8em]
\Large Полная согласованная версия:\\
от мереологического основания к реляционной физике}
\author{Andrey Tikhonov (Тихонов Андрей)\\\texttt{uwt@xteam.pro}}
\date{}

\begin{document}

\frontmatter
\maketitle

\chapter*{Аннотация}
\addcontentsline{toc}{chapter}{Аннотация}

Настоящий текст --- полная, внутренне согласованная и непротиворечивая редакция
Теории Единого Целого (ТЕЦ). Цель теории --- построить язык, в котором
пространство, время, движение, расстояние, масса и энергия не вводятся как
первичные понятия, а выводятся как производные структуры из более
фундаментального уровня: Единого Целого, его частей и отношений между ними.

Исходный принцип формулируется кратко:
\[
\boxed{\text{не отношения существуют в пространстве и времени,}}
\]
\[
\boxed{\text{а пространство и время возникают из отношений.}}
\]

В отличие от предыдущих версий, настоящая редакция:

\begin{itemize}
    \item строит основание не на теории множеств, а на мереологии (исчислении
    «часть"--"целое») над ограниченной решёткой, что устраняет скрытую
    пространственно-контейнерную интуицию;
    \item аксиоматизирует структуру значений отношений как нормированную
    упорядоченную группу с полным набором аксиом;
    \item \emph{выводит} метрику (неравенство треугольника становится теоремой),
    а не постулирует её;
    \item устраняет скрытую арифметику над значениями отношений, допуская числовые
    операции только после нормировки;
    \item \emph{выводит} кинетическую энергию $\tfrac12 m v^2$ как разложение
    функционала стоимости изменения, а не заимствует её;
    \item доказывает собственную непротиворечивость построением явной модели;
    \item извлекает проверяемые в принципе различающие предсказания.
\end{itemize}

\tableofcontents

\mainmatter
\setcounter{chapter}{-1}

%==================================================================
\chapter{Теория Абсолюта}
%==================================================================

\section{Абсолют как предельное основание}

Прежде аксиоматики необходимо зафиксировать концептуальный слой. Единое Целое не
мыслится как объект, расположенный в некотором внешнем пространстве или существующий
во внешнем времени. Напротив, пространство и время понимаются как внутренние
производные структуры самого Единого Целого.

Назовём этот уровень Теорией Абсолюта. Его задача --- установить, что Абсолют не
содержится в пространстве и времени, а является условием возможности их возникновения.

\begin{postulate}[Внепространственность Абсолюта]
Абсолютное Целое $\Whole$ не находится в пространстве:
\[
\Whole\notin\Space.
\]
Пространство, если оно возникает, возникает внутри $\Whole$ как структура отношений
его частей.
\end{postulate}

\begin{postulate}[Вневременность Абсолюта]
Абсолютное Целое $\Whole$ не существует во внешнем времени:
\[
\Whole\notin\Time.
\]
Время, если оно возникает, возникает внутри $\Whole$ как мера изменения состояний
отношений.
\end{postulate}

\begin{postulate}[Внутренность пространства и времени]
Пространство и время не содержат Абсолют. Они являются его внутренними производными
структурами:
\[
\boxed{\text{не Вселенная находится в пространстве и времени,}}
\]
\[
\boxed{\text{а пространство и время возникают внутри Вселенной.}}
\]
\end{postulate}

\section{Принцип возникновения различимости}

Более фундаментальным, чем отношение, является условие различимости. Отношение можно
задать лишь тогда, когда имеются различимые части. Без различимости невозможно
определить ни расстояние, ни направление, ни движение, ни время.

\begin{definition}[Различимость]
Различимостью называется возможность отличить одну часть от другой:
\[
\Disc(A_i,A_j)\iff A_i\neq A_j.
\]
\end{definition}

\begin{theorem}[Одна часть не порождает различимости]
Если существует только одна часть, то невозможно определить расстояние, направление,
размерность, движение и время.
\end{theorem}

\begin{proof}
Все перечисленные понятия требуют сравнения, а сравнение требует по крайней мере двух
различимых элементов. При единственной части $A_1$ не существует $A_j\neq A_1$,
поэтому различимость не возникает, а с нею --- ни отношение, ни расстояние, ни время.
\end{proof}

\section{Распределённый центр}

Если у Единого Целого нет внешнего пространства, у него не может быть и внешнего
геометрического центра. Центр понимается не как выделенная точка, а как способ
внутреннего описания.

\begin{postulate}[Распределённый центр]
Для всякой части $A_i$ центр её внутреннего описания совпадает с центром Единого
Целого:
\[
\Center(A_i)=\Center(\Whole).
\]
Привилегированной точки-центра не существует; любая часть может быть выбрана
локальным центром описания.
\end{postulate}

\section{Три принципа Теории Абсолюта}

\begin{enumerate}
    \item Абсолют существует вне пространства и времени.
    \item Пространство и время не содержат Абсолют, а возникают как внутренние
    отношения и изменения его частей.
    \item Любая часть описывает Абсолют из собственной точки зрения, поэтому каждый
    внутренний наблюдатель вправе считать себя центром описания, не нарушая единства
    Целого.
\end{enumerate}

Эти принципы делают последующую аксиоматику естественной: отношения, пространство и
время появляются как необходимые следствия отсутствия внешнего и возникновения
различимости внутри Единого Целого.

%==================================================================
\chapter{Сознание и материя как онтологические уровни}
%==================================================================

\section{Сознание как единое целое}

В предельном онтологическом смысле сознание рассматривается не как объект среди
объектов и не как процесс внутри пространства-времени, а как целостное основание,
относительно которого возможны проявление, различимость и наблюдение.

\begin{axiom}[Сознание]
Сознание есть единое целое. Оно не имеет частей, не локализовано, не материально,
не находится в пространстве и не находится во времени.
\[
\boxed{\text{сознание не является объектом пространства и времени.}}
\]
Следовательно, к сознанию как к основанию неприменимы понятия ``раньше'',
``позже'', ``здесь'' и ``там''.
\end{axiom}

Из этой аксиомы следует, что сознание не возникает и не исчезает во временном
смысле. Возникновение и исчезновение предполагают время, а время относится только
к уровню различимых состояний и изменений.

\section{Материя как уровень проявления}

Материя вводится не как независимое от основания начало и не как нечто, возникшее
после сознания во времени. Материя есть уровень проявления, на котором появляются
различимые части, отношения между ними, пространство, время, движение и наблюдаемые
объекты.

\begin{axiom}[Материя]
Материя существует как проявление внутри сознания. Это не означает, что сознание
когда-то во времени создало материю. Материальный мир существует как аспект
проявления основания.
\[
\boxed{\text{материя логически зависит от сознания, но не возникает после него во времени.}}
\]
\end{axiom}

Вопрос ``что было до возникновения материи?'' не имеет строгого смысла в данной
схеме, поскольку слово ``до'' уже предполагает время, а время появляется только
на материальном уровне как структура изменения отношений.

\section{Начало материи как логическая зависимость}

\begin{axiom}[Начало материи]
Начало материи не является событием. Оно выражает не временную последовательность,
а логическую зависимость материального уровня от более фундаментального основания.
\[
\boxed{\text{сознание является онтологическим основанием материи, а не её временной причиной.}}
\]
\end{axiom}

Иными словами, материя имеет начало в сознании только в онтологическом и логическом
смысле. Подобно тому как множество всех возможных точек является основанием любой
конкретной точки, но не возникает раньше неё во времени, сознание является
основанием материального уровня без отношения ``до'' и ``после''.

\section{Иерархия уровней}

Получается следующая иерархия:

\begin{enumerate}
    \item сознание как единое целое;
    \item отсутствие внешнего пространства и внешнего времени на уровне основания;
    \item материя как уровень проявления;
    \item пространство как отношение между различимыми частями;
    \item время как отношение изменений между состояниями;
    \item движение, причинность и наблюдаемые объекты как структуры материального
    уровня.
\end{enumerate}

Поэтому пространство, время и движение не являются фундаментальными сущностями.
Они являются свойствами и производными структурами материального уровня.

\section{Эпистемологическое следствие}

Если сознание лежит вне пространства и времени, то материальный эксперимент не
может обнаружить его как объект. Любой эксперимент уже использует материю,
пространство, время и причинность, то есть работает внутри более низкого уровня
иерархии.

\begin{proposition}[Граница материального исследования]
Материальная система может исследовать только материальные отношения. Она не может
выйти за пределы собственной области определения средствами, принадлежащими этой
же области.
\end{proposition}

Следовательно, сознание не может стать объектом исследования материи в том же
смысле, в каком материальный объект становится объектом эксперимента. Сознание
может быть непосредственно доступно только самому сознанию.

\begin{proposition}[Неполная объективируемость субъекта]
Субъект не может полностью превратить самого себя в объект, оставаясь при этом
тем же субъектом.
\end{proposition}

Это соответствует известным ограничениям самописания в логике и математике: система
не всегда способна полностью описать собственные основания изнутри средствами,
которые уже принадлежат самой системе.

\section{Связь с Теорией Единого Целого}

В терминах Теории Единого Целого сказанное можно зафиксировать так:

\begin{enumerate}
    \item Единое Целое является абсолютным основанием.
    \item Сознание можно рассматривать как аспект Единого Целого, в котором
    отсутствуют пространство и время.
    \item Материя представляет собой уровень отношений внутри Единого Целого.
    \item Пространство возникает как отношение между различимыми частями.
    \item Время возникает как отношение изменений между состояниями этих частей.
    \item Пока нет различимых частей, нет ни пространства, ни времени.
\end{enumerate}

Ключевая формулировка этой главы:
\[
\boxed{\text{материя имеет основание в сознании, но не следует за ним во времени.}}
\]

%==================================================================
\chapter{Введение и порядок вывода}
%==================================================================

\section{Методологическое требование}

В теории запрещается использовать пространство, время, координаты, расстояние,
скорость, массу и энергию в качестве исходных понятий. Эти величины могут появиться
только после того, как определены отношения и состояния:
\[
\Rel\prec\Space,\quad \Rel\prec\Time,\quad \Rel\prec\Physics,
\]
где $\prec$ означает концептуальную первичность.

\section{Полная цепочка вывода}

Фундаментальная последовательность настоящей версии имеет вид
\[
(\,\mathfrak U,\PartOf,\Whole\,)
\to \Parts
\to \Disc
\to \Rel
\to \Vstruct
\to \State
\to \Time
\to d
\to v,a
\to m,p,\mathfrak F,E,L,H.
\]
Каждая стрелка в дальнейшем снабжается либо аксиомой, либо доказанной теоремой. В
этом состоит требование согласованности: ни одно производное понятие не вводится
раньше, чем определены его предпосылки.

%==================================================================
\chapter{Основание: мереология и решётка}
%==================================================================

\section{Почему не теория множеств}

Запись $A\subseteq\Whole$, $X\in\Parts$ опирается на теоретико-множественную
интуицию, в которой множество мыслится как вместилище элементов и несёт скрытую
пространственно-контейнерную семантику. Это противоречит цели ТЕЦ. Поэтому основание
строится на мереологии --- исчислении отношения «часть"--"целое» $\PartOf$ --- над
ограниченным частичным порядком. Отношение $\PartOf$ фиксирует лишь структурный
порядок: ни точек, ни вместимости, ни расположения.

\section{Мереологические аксиомы}

Единственный первичный нелогический символ --- бинарное отношение $\PartOf$ на области
рассуждения $\mathfrak U$.

\begin{axiom}[Рефлексивность]
$x\PartOf x$.
\end{axiom}

\begin{axiom}[Антисимметрия]
Если $x\PartOf y$ и $y\PartOf x$, то $x=y$.
\end{axiom}

\begin{axiom}[Транзитивность]
Если $x\PartOf y$ и $y\PartOf z$, то $x\PartOf z$.
\end{axiom}

Тем самым $\PartOf$ есть частичный порядок.

\begin{axiom}[Единое Целое как наибольший элемент]
Существует единственный $\Whole$ с $x\PartOf\Whole$ для всех $x$:
\[
\exists!\,\Whole\ \forall x:\ x\PartOf\Whole.
\]
\end{axiom}

\begin{axiom}[Реляционная полнота]
Для всякого непустого совокупления частей существует точная верхняя грань. Структура
$(\mathfrak U,\PartOf,\Whole)$ есть ограниченная полная верхняя полурешётка.
\end{axiom}

\begin{definition}[Часть и класс частей]
Собственная часть:
\[
x\PropPart\Whole\;:\Longleftrightarrow\;x\PartOf\Whole\ \wedge\ x\neq\Whole.
\]
Класс частей: $\Parts:=\{x\mid x\PropPart\Whole\}$.
\end{definition}

\begin{definition}[Различимость]
$\Disc(x,y):\Longleftrightarrow x\neq y$.
\end{definition}

\begin{proposition}[Отсутствие контейнерной семантики]
В основании $(\mathfrak U,\PartOf,\Whole)$ нет ни принадлежности $\in$, ни точек, ни
вместимости: все отношения суть отношения порядка.
\end{proposition}

\begin{proof}
Сигнатура содержит единственный нелогический символ $\PartOf$ --- частичный порядок.
Понятия точки, расположения внутри и объёма не определимы средствами одного лишь
порядка; они появляются лишь после введения отношений $\Rel$ и нормировки $\norm$ как
производные.
\end{proof}

\begin{remark}[Метатеория]
Чтобы и понятие отображения не вносило теоретико-множественной интуиции, метатеорию
допустимо принять типовой: «отображение» понимается как терм функционального типа, а
не как множество пар.
\end{remark}

%==================================================================
\chapter{Структура значений отношений}
%==================================================================

\section{Нормированная упорядоченная группа значений}

Значения отношений должны сравниваться, компоноваться и проецироваться в число. Все
эти операции аксиоматически замкнуты, а единственным мостом к числам служит
нормировка $\norm$.

\begin{definition}[Структура значений]
Структурой значений отношений называется набор
\[
\Vstruct=(\Values,\comp,\ide,\Vord,\norm).
\]
\end{definition}

\begin{axiom}[Группа]
$(\Values,\comp,\ide)$ есть группа: $\comp$ ассоциативна, $\ide$ нейтрален, каждый
$a$ обратим:
\[
a\comp\ide=\ide\comp a=a,\qquad a\comp a^{-1}=a^{-1}\comp a=\ide,\qquad
(a\comp b)^{-1}=b^{-1}\comp a^{-1}.
\]
\end{axiom}

\begin{axiom}[Порядок]
$\Vord$ есть частичный порядок на $\Values$.
\end{axiom}

\begin{axiom}[Монотонность композиции]
\[
a\Vord b\ \Longrightarrow\ a\comp c\Vord b\comp c\ \wedge\ c\comp a\Vord c\comp b.
\]
\end{axiom}

\begin{axiom}[Нормировка]\label{ax:norm}
Существует $\norm:\Values\to\mathbb R_{\ge0}$ со свойствами
\begin{align}
&\text{(N1)}\quad \norm(\ide)=0; \notag\\
&\text{(N2)}\quad \norm(a)=0\ \Longrightarrow\ a=\ide \quad\text{(положительная определённость)};\notag\\
&\text{(N3)}\quad \norm(a^{-1})=\norm(a)\quad\text{(симметрия)};\notag\\
&\text{(N4)}\quad \norm(a\comp b)\le\norm(a)+\norm(b)\quad\text{(субаддитивность)};\notag\\
&\text{(N5)}\quad a\Vord b\ \Longrightarrow\ \norm(a)\le\norm(b)\quad\text{(монотонность)}.\notag
\end{align}
\end{axiom}

\begin{remark}[О согласованности аксиом значений]
В настоящей редакции \emph{не} постулируется, что $\ide$ --- наименьший элемент
порядка. Такое требование в сочетании с обратимостью группы привело бы к
противоречию: из $\ide\Vord a$ и $\ide\Vord a^{-1}$ по монотонности следовало бы
$a\Vord\ide$ и потому $a=\ide$, то есть группа стала бы тривиальной. Отказ от этой
избыточной аксиомы --- необходимое условие непротиворечивости (см. главу~12).
\end{remark}

\begin{proposition}[$\norm$ есть абстрактная длина]
$\norm(a)\ge0$, $\;\norm(a)=0\iff a=\ide$, $\;\norm(a^{-1})=\norm(a)$,
$\;\norm(a\comp b)\le\norm(a)+\norm(b)$.
\end{proposition}

\begin{proof}
Неотрицательность --- по области значений. Эквивалентность $\norm(a)=0\iff a=\ide$:
$\Leftarrow$ есть N1, $\Rightarrow$ есть N2. Симметрия --- N3, субаддитивность --- N4.
\end{proof}

%==================================================================
\chapter{Отношения}
%==================================================================

\section{Отношение как отображение в группу значений}

\begin{definition}[Отношение]
Отношением называется отображение
\[
\Rel:\Parts\times\Parts\to\Values.
\]
\end{definition}

\begin{axiom}[Рефлексивность отношения]
$\Rel(A,A)=\ide$.
\end{axiom}

\begin{axiom}[Различимость и нетривиальность]\label{ax:nontrivial}
\[
\Rel(A_i,A_j)=\ide\ \Longleftrightarrow\ A_i=A_j.
\]
\end{axiom}

\begin{axiom}[Обратимость отношения]\label{ax:invol}
\[
\Rel(A_j,A_i)=\Rel(A_i,A_j)^{-1}.
\]
\end{axiom}

\begin{axiom}[Композиция отношений]\label{ax:compose}
Для любых частей $A_i,A_j,A_k$:
\[
\Rel(A_i,A_k)\ \Vord\ \Rel(A_i,A_j)\comp\Rel(A_j,A_k).
\]
\end{axiom}

\begin{remark}
Аксиома~\ref{ax:compose} полностью внутренняя: в ней нет ни чисел, ни расстояний.
Содержательно она выражает принцип «прямое отношение не слабее составного»: посредник
не делает связь теснее прямой. Именно эта аксиома станет источником неравенства
треугольника.
\end{remark}

\begin{proposition}[Совместность аксиом отношений]
Аксиомы~\ref{ax:nontrivial}--\ref{ax:compose} совместны при подстановке совпадающих
индексов: при $j=i$ и при $k=i$ аксиома~\ref{ax:compose} обращается в рефлексивное
$\Rel(A_i,A_k)\Vord\Rel(A_i,A_k)$ и $\ide\Vord\ide$ соответственно.
\end{proposition}

\begin{proof}
При $j=i$: $\Rel(A_i,A_i)\comp\Rel(A_i,A_k)=\ide\comp\Rel(A_i,A_k)=\Rel(A_i,A_k)$.
При $k=i$: $\Rel(A_i,A_j)\comp\Rel(A_j,A_i)=\Rel(A_i,A_j)\comp\Rel(A_i,A_j)^{-1}=\ide$,
а $\Rel(A_i,A_i)=\ide$. Оба случая дают рефлексивное неравенство.
\end{proof}

%==================================================================
\chapter{Состояния и изменение}
%==================================================================

\section{Системы и состояния}

\begin{definition}[Система частей]
Системой частей называется непустое подмножество $\SubParts\subseteq\Parts$.
\end{definition}

\begin{definition}[Состояние]
Состоянием системы $\SubParts$ называется совокупность отношений между её различными
частями:
\[
\State(\SubParts):=\{\Rel(A_i,A_j)\mid A_i,A_j\in\SubParts,\ i\neq j\}.
\]
\end{definition}

\begin{definition}[Изменение]
Изменением называется упорядоченная пара различных состояний
$\Delta\State=(\State_1,\State_2)$, $\State_1\neq\State_2$.
\end{definition}

\section{Реляционное приращение}

\begin{definition}[Приращение отношения]\label{def:increment}
Пусть $\Rel_k(A_i,A_j)$ и $\Rel_{k+1}(A_i,A_j)$ --- значения отношения в состояниях
$\State_k,\State_{k+1}$. Приращением называется элемент группы
\[
\delta_{ij}:=\Rel_k(A_i,A_j)^{-1}\comp\Rel_{k+1}(A_i,A_j),
\]
единственный с $\Rel_{k+1}(A_i,A_j)=\Rel_k(A_i,A_j)\comp\delta_{ij}$.
\end{definition}

\begin{remark}[Чистота операций]
Приращение определено лишь внутренними операциями группы значений (композиция и
обращение), без какой-либо арифметики над $\Values$.
\end{remark}

%==================================================================
\chapter{Выведение пространства: метрика как теорема}
%==================================================================

\section{Реляционное расстояние}

\begin{definition}[Реляционное расстояние]
\[
d(A_i,A_j):=\norm\bigl(\Rel(A_i,A_j)\bigr).
\]
\end{definition}

\begin{theorem}[Метрика выводится из аксиом]\label{thm:metric}
Отображение $d:\Parts\times\Parts\to\mathbb R_{\ge0}$ есть метрика; все четыре её
аксиомы суть теоремы.
\end{theorem}

\begin{proof}
\textbf{Неотрицательность.} $d=\norm(\cdot)\ge0$.

\textbf{Тождество неразличимых.} $\norm(a)=0\iff a=\ide$, и по
аксиоме~\ref{ax:nontrivial} $\Rel(A_i,A_j)=\ide\iff A_i=A_j$; значит $d(A_i,A_j)=0\iff
A_i=A_j$.

\textbf{Симметрия.} По обратимости отношения и симметрии нормы:
\[
d(A_j,A_i)=\norm(\Rel(A_i,A_j)^{-1})=\norm(\Rel(A_i,A_j))=d(A_i,A_j).
\]

\textbf{Неравенство треугольника.} По
аксиоме~\ref{ax:compose}, монотонности (N5) и субаддитивности (N4):
\[
d(A_i,A_k)=\norm(\Rel(A_i,A_k))
\le\norm(\Rel(A_i,A_j)\comp\Rel(A_j,A_k))
\le\norm(\Rel(A_i,A_j))+\norm(\Rel(A_j,A_k))
=d(A_i,A_j)+d(A_j,A_k).
\]
Все четыре свойства установлены.
\end{proof}

\begin{definition}[Реляционное пространство]
Реляционным пространством системы $\SubParts$ называется метрическое пространство
$\Space(\SubParts):=(\SubParts,d)$, порождённое нормировкой отношений.
\end{definition}

\begin{corollary}[Критерий метризуемости]
Реляционная система метризуема тогда и только тогда, когда значения её отношений
образуют нормированную группу $\Vstruct$ с аксиомой композиции. Привычное
пространство --- частный случай реляционной структуры, а не её необходимое свойство.
\end{corollary}

\section{Геометрия как инвариант}

\begin{definition}[Автоморфизм]
Автоморфизмом реляционной структуры $\SubParts$ называется биекция
$\phi:\SubParts\to\SubParts$, сохраняющая отношения:
$\Rel(A_i,A_j)=\Rel(\phi(A_i),\phi(A_j))$.
\end{definition}

\begin{definition}[Геометрический инвариант]
Геометрическим инвариантом называется величина, сохраняющаяся при всех автоморфизмах.
Геометрия системы определяется классом её инвариантов.
\end{definition}

\begin{proposition}
Всякий автоморфизм есть изометрия метрики $d$.
\end{proposition}

\begin{proof}
$d(\phi(A_i),\phi(A_j))=\norm(\Rel(\phi(A_i),\phi(A_j)))=\norm(\Rel(A_i,A_j))=d(A_i,A_j)$.
\end{proof}

%==================================================================
\chapter{Выведение времени}
%==================================================================

\section{Время как структура изменения}

В базовой версии время сводилось к мере различия состояний. Этого недостаточно:
необходимо различать длительность, порядок и необратимость.

\begin{definition}[Длительность]
Длительностью называется отображение $\tau:\States\times\States\to\mathbb R_{\ge0}$
с
\[
\tau(\State_i,\State_j)=0\iff\State_i=\State_j,\qquad
\tau(\State_i,\State_j)=\tau(\State_j,\State_i).
\]
\end{definition}

\begin{definition}[Направление]
Направление времени задаётся отношением порядка $\prec$ на состояниях: запись
$\State_i\prec\State_j$ означает, что $\State_j$ следует после $\State_i$ в данной
истории.
\end{definition}

\begin{definition}[Необратимость]
Необратимость задаётся функционалом $\Omega:\States\to\mathbb R$ (сложность,
энтропия), для которого вдоль истории $\State_0\prec\State_1\prec\cdots$ выполняется
$\Omega(\State_{k+1})\ge\Omega(\State_k)$.
\end{definition}

\begin{definition}[Реляционное время]
Реляционным временем называется тройка
\[
\Time=(\tau,\prec,\Omega).
\]
\end{definition}

\begin{theorem}[Атемпоральность статической системы]
Если история постоянна, $\State_0=\State_1=\cdots$, то время не возникает:
$\tau(\State_k,\State_{k+1})=0$ для всех $k$.
\end{theorem}

\begin{proof}
По определению длительности $\tau(\State,\State)=0$. Отсутствие различимого изменения
влечёт отсутствие времени.
\end{proof}

%==================================================================
\chapter{Наблюдатели и относительные центры}
%==================================================================

\begin{definition}[Внутренний наблюдатель]
Внутренним наблюдателем называется часть $O\in\Parts$, относительно которой строится
описание:
\[
\Observer_O:=\{\Rel(O,A_j)\mid A_j\in\Parts,\ A_j\neq O\}.
\]
\end{definition}

\begin{theorem}[Каждая часть может быть центром описания]
Для любой части $A_i$ можно построить описание системы относительно неё.
\end{theorem}

\begin{proof}
Совокупность $\mathcal R_{A_i}=\{\Rel(A_i,A_j)\mid A_j\neq A_i\}$ задаёт описание
относительно $A_i$; так как $A_i$ произвольна, описание возможно для любой части.
\end{proof}

\begin{definition}[Согласованность наблюдателей]
Описания $\mathcal R_{A_i}$ и $\mathcal R_{A_k}$ согласованы, если существует
преобразование $\Phi_{ik}$, сохраняющее инварианты структуры:
$\Inv(\mathcal R_{A_i})=\Inv(\Phi_{ik}(\mathcal R_{A_i}))$.
\end{definition}

\begin{proposition}[Объективность инвариантов]
Геометрические инварианты не зависят от выбора наблюдателя.
\end{proposition}

\begin{proof}
Согласующее преобразование $\Phi_{ik}$ по определению сохраняет $\Inv$; значит
значение инварианта одинаково во всех согласованных описаниях.
\end{proof}

%==================================================================
\chapter{Реляционная механика}
%==================================================================

\section{Движение, скорость, ускорение}

\begin{definition}[Движение]
Движением части $A_i$ относительно $A_j$ называется изменение отношения:
\[
\Motion(A_i,A_j):=\delta_{ij}\ \text{(см. опр.~\ref{def:increment})}.
\]
Движение не есть перемещение в заранее данном пространстве; движение есть изменение
отношения.
\end{definition}

\begin{definition}[Скорость]
\[
v(A_i,A_j):=\frac{\Delta d(A_i,A_j)}{\Delta t},\qquad
\Delta d=\norm(\delta_{ij}),\quad \Delta t=\tau(\State_k,\State_{k+1}).
\]
\end{definition}

\begin{definition}[Ускорение]
$\displaystyle a(A_i,A_j):=\frac{\Delta v(A_i,A_j)}{\Delta t}$.
\end{definition}

\section{Устойчивость и масса без скрытой арифметики}

\begin{definition}[Вариация части]
Вариацией части $A_i$ называется число, получаемое \emph{после} нормировки каждого
приращения:
\[
\Varr(A_i;\State_k,\State_{k+1}):=\sum_{j\neq i}\norm(\delta_{ij})\in\mathbb R_{\ge0}.
\]
\end{definition}

\begin{definition}[Устойчивость и масса]
\[
\Stab(A_i):=\frac{1}{1+\Varr(A_i)},\qquad
m(A_i):=M\bigl(\Stab(A_i)\bigr),
\]
где $M:[0,1]\to\mathbb R_{\ge0}$ --- фиксированная монотонная калибровочная функция.
\end{definition}

\begin{proposition}[Чистота арифметики]
В определениях вариации, устойчивости и массы нет ни одной арифметической операции над
значениями $\Values$: единственный переход к $\mathbb R$ осуществляется нормировкой
$\norm$, после чего все операции выполняются над числами.
\end{proposition}

\begin{proof}
Приращение $\delta_{ij}$ построено внутренними групповыми операциями. Норма
$\norm(\delta_{ij})\in\mathbb R_{\ge0}$. Сумма, прибавление единицы и деление
действуют лишь над этими числами.
\end{proof}

\section{Импульс и сила}

\begin{definition}[Импульс]
$p_{ij}:=m(A_i)\,v(A_i,A_j)$.
\end{definition}

\begin{definition}[Сила]
$\displaystyle\mathfrak F_{ij}:=\frac{\Delta p_{ij}}{\Delta t}$.
\end{definition}

%==================================================================
\chapter{Энергия: вывод кинетической формы}
%==================================================================

\section{Функционал стоимости изменения}

Множитель $\tfrac12$ и квадрат скорости не заимствуются, а выводятся как структура
разложения функционала стоимости изменения вблизи равновесия.

\begin{definition}[Функционал стоимости]
Функционалом стоимости называется $\Cost:\mathbb R\to\mathbb R_{\ge0}$,
сопоставляющий реляционной скорости $v$ внутреннюю стоимость поддержания изменения.
\end{definition}

\begin{postulate}[Свойства стоимости]\label{post:cost}
\begin{enumerate}
    \item[(E1)] $\Cost(0)=0$ (покой ничего не стоит);
    \item[(E2)] минимум при $v=0$ (покой устойчив), откуда $\Cost'(0)=0$;
    \item[(E3)] $\Cost\in C^2$ в окрестности нуля;
    \item[(E4)] чётность $\Cost(v)=\Cost(-v)$ (изотропия), откуда все нечётные
    производные в нуле равны нулю.
\end{enumerate}
\end{postulate}

\begin{theorem}[Квадратичность кинетической энергии]\label{thm:kinetic}
При E1--E4 для малых скоростей
\[
\Cost(v)=\tfrac12\,m\,v^2+O(v^4),\qquad m:=\Cost''(0)\ge0.
\]
Множитель $\tfrac12$ есть тейлоровский коэффициент $1/2!$, а масса $m$ --- кривизна
функционала стоимости в точке равновесия.
\end{theorem}

\begin{proof}
Разложение Тейлора в нуле:
\[
\Cost(v)=\Cost(0)+\Cost'(0)v+\tfrac12\Cost''(0)v^2+\tfrac16\Cost'''(0)v^3+O(v^4).
\]
По E1 $\Cost(0)=0$; по E2 $\Cost'(0)=0$; по E4 $\Cost'''(0)=0$. Значит
$\Cost(v)=\tfrac12\Cost''(0)v^2+O(v^4)$. Полагаем $m:=\Cost''(0)$; в точке минимума
$\Cost''(0)\ge0$.
\end{proof}

\begin{corollary}[Согласование двух определений массы]
$m=\Cost''(0)$ согласовано с массой через устойчивость: большая кривизна означает
крутой рост стоимости изменения, то есть сильное сопротивление изменению отношений,
то есть высокую устойчивость $\Stab(A_i)$. Инерция как кривизна и инерция как
устойчивость измеряют одно и то же.
\end{corollary}

\section{Полная энергия, лагранжиан и гамильтониан}

\begin{definition}[Кинетическая и потенциальная части]
\[
K(\SubParts)=\sum_{i,j}\tfrac12 m(A_i)v_{ij}^2,\qquad
P(\SubParts)=\sum_{i,j}\Phi\bigl(\Rel(A_i,A_j)\bigr),
\]
где $\Phi$ --- потенциал отношений.
\end{definition}

\begin{definition}[Энергия, лагранжиан, гамильтониан]
\[
E=K+P,\qquad L=K-P,\qquad H=K+P.
\]
\end{definition}

\begin{remark}
Кинетическая часть имеет квадратичную форму не по соглашению, а по
теореме~\ref{thm:kinetic}; следующий по порядку член $O(v^4)$ порождает проверяемые
поправки (глава~13).
\end{remark}

%==================================================================
\chapter{Непротиворечивость теории}
%==================================================================

\section{Постановка вопроса}

Теория согласована, если её аксиомы совместны, то есть существует интерпретация
(модель), в которой все они истинны. Построим такую модель; тем самым будет доказана
\emph{относительная} непротиворечивость: если непротиворечива евклидова геометрия
(а значит арифметика вещественных чисел), то непротиворечива и ТЕЦ.

\section{Евклидова модель}

\begin{definition}[Стандартная модель $\mathfrak M$]\label{def:model}
Зафиксируем размерность $n\ge1$ и инъективное сопоставление частям точек:
\[
A_i\longmapsto x_i\in\mathbb R^n,\qquad A_i\neq A_j\Rightarrow x_i\neq x_j.
\]
Положим:
\begin{itemize}
    \item группа значений $\Values=(\mathbb R^n,+,0)$, обращение $a^{-1}=-a$;
    \item порядок $\Vord$ --- дискретный: $a\Vord b\iff a=b$;
    \item нормировка $\norm(a)=\|a\|$ (евклидова норма);
    \item отношение $\Rel(A_i,A_j):=x_j-x_i$;
    \item длительность $\tau(\State_k,\State_{k+1})$ --- любая положительная величина
    различия состояний; стоимость $\Cost(v)=\tfrac12 v^2$.
\end{itemize}
\end{definition}

\begin{theorem}[$\mathfrak M$ удовлетворяет всем аксиомам]\label{thm:consistency}
Модель $\mathfrak M$ из определения~\ref{def:model} удовлетворяет всем аксиомам
структуры значений (глава~4) и всем аксиомам отношений (глава~5).
\end{theorem}

\begin{proof}
\emph{Группа.} $(\mathbb R^n,+,0)$ --- абелева группа; ассоциативность,
нейтральность $0$ и обратимость $-a$ выполнены.

\emph{Порядок и монотонность.} Дискретный порядок ($a\Vord b\iff a=b$) есть частичный
порядок; монотонность тривиальна: $a=b\Rightarrow a+c=b+c$.

\emph{Нормировка.} N1: $\|0\|=0$. N2: $\|a\|=0\Rightarrow a=0$. N3:
$\|-a\|=\|a\|$. N4: $\|a+b\|\le\|a\|+\|b\|$ --- неравенство треугольника евклидовой
нормы. N5: $a=b\Rightarrow\|a\|=\|b\|$.

\emph{Рефлексивность отношения.} $\Rel(A,A)=x-x=0=\ide$.

\emph{Различимость и нетривиальность.} $\Rel(A_i,A_j)=0\iff x_i=x_j\iff A_i=A_j$
(инъективность сопоставления).

\emph{Обратимость.} $\Rel(A_j,A_i)=x_i-x_j=-(x_j-x_i)=\Rel(A_i,A_j)^{-1}$.

\emph{Композиция.} $\Rel(A_i,A_j)\comp\Rel(A_j,A_k)=(x_j-x_i)+(x_k-x_j)=x_k-x_i
=\Rel(A_i,A_k)$, поэтому $\Rel(A_i,A_k)\Vord\Rel(A_i,A_j)\comp\Rel(A_j,A_k)$ выполнено
с равенством (дискретный порядок рефлексивен).

Все аксиомы истинны в $\mathfrak M$.
\end{proof}

\begin{corollary}[Относительная непротиворечивость ТЕЦ]
Система аксиом ТЕЦ непротиворечива при условии непротиворечивости арифметики
вещественных чисел.
\end{corollary}

\begin{proof}
Если бы из аксиом ТЕЦ выводилось противоречие, оно было бы истинно и в модели
$\mathfrak M$, построенной средствами $\mathbb R^n$. Это означало бы противоречивость
арифметики $\mathbb R$. Следовательно, при непротиворечивости последней непротиворечива
и ТЕЦ.
\end{proof}

\begin{theorem}[Согласие модели с выведенными понятиями]
В модели $\mathfrak M$ реляционное расстояние совпадает с евклидовым, а кинетическая
энергия --- с ньютоновской:
\[
d(A_i,A_j)=\|x_j-x_i\|,\qquad \Cost(v)=\tfrac12 v^2.
\]
\end{theorem}

\begin{proof}
$d(A_i,A_j)=\norm(\Rel(A_i,A_j))=\|x_j-x_i\|$ по определению. Стоимость
$\Cost(v)=\tfrac12 v^2$ удовлетворяет E1--E4 с $\Cost''(0)=1$, что согласуется с
теоремой~\ref{thm:kinetic} при $m=1$.
\end{proof}

\begin{remark}[Содержательность теории]
Модель $\mathfrak M$ плоская (порядок дискретен, композиция точна). Аксиомы, однако,
допускают и неплоские реализации, где аксиома~\ref{ax:compose} выполняется со строгим
неравенством, а порядок $\Vord$ нетривиален. Тем самым теория непротиворечива и при
этом не сводится к одной лишь евклидовой геометрии.
\end{remark}

%==================================================================
\chapter{Физическая интерпретация и предсказания}
%==================================================================

\section{Сопоставление со схемами}

Классическая схема: $\text{пространство-время}\to\text{материя}\to\text{движение}$.
Схема ТЕЦ:
\[
\Whole\to\Parts\to\Disc\to\Rel\to\State\to\Space,\Time\to\text{физические величины}.
\]
Пространство и время не сцена, а устойчивые формы организации отношений.

\section{Различающие предсказания}

\begin{proposition}[A. Неньютоновские поправки]
Из теоремы~\ref{thm:kinetic} кинетическая энергия имеет вид
$K(v)=\tfrac12 mv^2+c_4v^4+O(v^6)$, $c_4=\tfrac1{24}\Cost^{(4)}(0)$. ТЕЦ предсказывает
отклонение от ньютоновской зависимости фиксированной (чётной) формы.
\end{proposition}

\begin{theorem}[B. Инвариантная предельная скорость]
Если $\sup_a\norm(a)=\norm_{\max}<\infty$ и длительность имеет минимальный шаг
$\tau_{\min}>0$, то реляционная скорость ограничена инвариантой
\[
v\le\frac{\norm_{\max}}{\tau_{\min}}=:c.
\]
\end{theorem}

\begin{proof}
$\Delta d=\norm(\delta)\le\norm_{\max}$, $\Delta t\ge\tau_{\min}$, откуда
$v=\Delta d/\Delta t\le\norm_{\max}/\tau_{\min}$. Константы структурны и одинаковы для
всех частей, поэтому $c$ инвариантна.
\end{proof}

\begin{theorem}[C. Квантование длины и времени]
Если порядок $(\Values,\Vord)$ фундирован и содержит атомы $a_0$ с
$\norm(a_0)=\ell_0>0$, то ненулевые расстояния не меньше $\ell_0$, а длительность
имеет минимальный шаг $\tau_0>0$.
\end{theorem}

\begin{proof}
Любое $a\neq\ide$ мажорирует атом $a_0\Vord a$; по монотонности
$\norm(a)\ge\norm(a_0)=\ell_0$. Аналогично для $\tau$.
\end{proof}

\begin{remark}[D. Размерная редукция]
Размерность есть размерность вложения метрики $(\Parts,d)$, то есть число независимых
инвариантов. На наименьших масштабах, где число различимых частей мало, это число
падает; ТЕЦ качественно предсказывает уменьшение реляционной размерности на малых
масштабах.
\end{remark}

\section{Критерии физической применимости}

\begin{enumerate}
    \item из отношений строится метрика (выполнено: теорема~\ref{thm:metric});
    \item из изменения состояний получается устойчивая временная параметризация
    (выполнено: глава~8);
    \item определены наблюдатели и преобразования между описаниями (глава~9);
    \item физические величины выражены через реляционные структуры (главы~10--11);
    \item теория даёт проверяемые отличия (предсказания A--C).
\end{enumerate}

%==================================================================
\chapter{Заключение}
%==================================================================

Полная версия Теории Единого Целого реализует исходный замысел без внутренних
противоречий. Главные результаты:

\begin{enumerate}
    \item основание построено на мереологии, а не на теории множеств, что устраняет
    скрытую пространственную интуицию;
    \item структура значений отношений аксиоматизирована как нормированная
    упорядоченная группа, причём набор аксиом проверен на совместность;
    \item метрика, включая неравенство треугольника, \emph{выведена} как теорема;
    \item арифметика над значениями устранена: число появляется только через
    нормировку;
    \item кинетическая энергия $\tfrac12 mv^2$ \emph{выведена} как разложение
    функционала стоимости, а масса отождествлена с его кривизной;
    \item непротиворечивость доказана построением явной евклидовой модели;
    \item получены три различающих предсказания и одно качественное.
\end{enumerate}

В краткой форме принцип теории:
\[
\boxed{\Rel\prec\Space,\Time,\Physics,}
\]
\[
\boxed{\text{метрика, масса и энергия выводятся из нормированной композиции отношений.}}
\]

\backmatter

\chapter*{Сводка решённых проблем}
\addcontentsline{toc}{chapter}{Сводка решённых проблем}

\begin{center}
\renewcommand{\arraystretch}{1.4}
\begin{tabular}{>{\raggedright\arraybackslash}p{0.30\textwidth}>{\raggedright\arraybackslash}p{0.62\textwidth}}
\toprule
\textbf{Проблема прежних версий} & \textbf{Решение в полной версии} \\
\midrule
Круговость метрики & Неравенство треугольника доказано (теорема~\ref{thm:metric}) из
аксиомы композиции~\ref{ax:compose}. \\
Скрытая арифметика & Приращение определено в группе значений; $\norm$ применяется до
арифметики. \\
Заимствованная энергия & $\tfrac12 mv^2$ выведено (теорема~\ref{thm:kinetic}); масса
$=\Cost''(0)$. \\
Множественное основание & Мереология над решёткой $(\mathfrak U,\PartOf,\Whole)$;
без $\in$ и точек. \\
Незамкнутая структура $\Values$ & Полная аксиоматизация нормированной упорядоченной
группы; совместность проверена. \\
Отсутствие предсказаний & Поправки $O(v^4)$, предельная скорость, квантование,
размерная редукция. \\
Недоказанная непротиворечивость & Построена евклидова модель
(теорема~\ref{thm:consistency}); доказана относительная непротиворечивость. \\
\bottomrule
\end{tabular}
\end{center}

\end{document}
